\section{Conclusion}

Public benchmark rank and deployment fit can be genuinely different questions, and in our
two-policy case study they were not even the same \emph{kind} of question: the two
policies' public scores were not commensurable at all. We proposed GFIC, a four-axis
diagnostic protocol anchored to independent causal-inference and interpretability
literatures, and showed on a ghost-probe target scenario that it localizes
benchmark-to-deployment gaps to specific, actionable links of a policy's causal chain --
a repair site for one policy that differs from another's, a causal-confusion signature
invisible to open-loop scoring, and a domain-robustness ordering that reverses between
representation and behaviour. We further showed, honestly, that repairing a diagnosed
link does not guarantee a behavioural gain at a small budget, a finding that survived
only because of an attribution-control arm designed to rule out exactly this confound.
We deliberately stopped short of a single composite score, and short of a fully closed
three-way comparison against post-training outcomes at scale; both are natural next
steps, together with extending the protocol to additional scenario types and to policies
whose native training domain can be independently verified.
